% P17-D2 continuation: finite-state bets and subjective probability.
\chapter{Finite Bets and Subjective Probabilities}
\label{sec:p17-finite-subjective-probability}

In a lottery model the probabilities are part of the specified
alternative. A state-contingent bet instead lists what is paid in each
possible state. Its probability weights can be derived from a preference
only under further assumptions. This chapter gives Rubinstein's finite
betting presentation of the de Finetti approach
\cite[2015 revision, pp.~102--103]{Rubinstein2015LectureNotes}.
It proves the representation completely and then derives a finite
fair-price coherence theorem as an explicit handbook consequence.

The representation below averages monetary payoffs themselves. Its
additivity assumption is substantially stronger than the mere existence
of beliefs about states. A separate example demonstrates that distinction.
No observed frequency or uniquely correct belief is inferred from the
representation theorem.

\section{State-contingent payoffs and preference assumptions}
\label{sec:p17-finite-bet-domain}

\begin{hbdefinition}{def:p17-finite-bets}{States, bets and certain payments}
Let \(S=\{1,\ldots,K\}\), where \(K\geq1\) is an integer, be a
finite exhaustive set of mutually exclusive states. A \emph{bet} is a
vector \(x=(x_1,\ldots,x_K)\in\mathbb R^K\); it pays the real amount
\(x_i\) in state \(i\). Negative amounts are losses in the same
payoff units and at the same settlement time.
Vector addition and scalar multiplication operate coordinatewise.

Write \(\mathbf0=(0,\ldots,0)\), \(\mathbf1=(1,\ldots,1)\),
and let \(e_i\) be the vector with coordinate one at \(i\) and zero
elsewhere. The vector \(t\mathbf1\) is a sure payment of \(t\).
Write \(x\geq_{\rm cw}y\) for \(x_i\geq y_i\) at every coordinate.
The usual Euclidean distance is
\(d(x,y)=(\sum_i(x_i-y_i)^2)^{1/2}\); also put
\(\|x\|_\infty=\max_i|x_i|\).

A \emph{weak preference} \(\succeq\) on all of \(\mathbb R^K\)
is complete and transitive. Its strict and indifferent comparisons,
\(\succ\) and \(\sim\), have their usual definitions by the two
weak comparisons. This is the source's bet space, with all finite
real payoff vectors permitted \cite[p.~102]{Rubinstein2015LectureNotes}.
\index{bet!state contingent}\index{subjective probability!finite states}
\end{hbdefinition}

\begin{hbdefinition}{def:p17-bet-axioms}{Continuity, strict improvement in every state, and additivity}
The following three conditions are imposed on the weak preference:
\begin{enumerate}
\item \emph{Continuity:} if \(x\succ y\), there is \(\varepsilon>0\)
such that \(x'\succ y'\) whenever
\(d(x',x)<\varepsilon\) and \(d(y',y)<\varepsilon\).
\item \emph{Weak monotonicity in the source's terminology:} if
\(x_i>y_i\) for every \(i\), then \(x\succ y\).
\item \emph{Additivity:} if \(x\succeq y\), then
\(x+z\succeq y+z\) for every \(z\in\mathbb R^K\).
\end{enumerate}
These are the source's three conditions
\cite[pp.~102--103]{Rubinstein2015LectureNotes}.
The second requires strict improvement in \emph{all} states. It does
not require strict preference from improvement in just one state.
The third permits an arbitrary common state-dependent payoff vector,
not merely addition of the same sure amount in every state.
\end{hbdefinition}

\begin{hblemma}{thm:p17-bet-cash-equivalent}{proved=annex:p17.bet-cash-equivalent}{The cash equivalent forced by the three axioms}
Under these assumptions, adding a common vector preserves weak, strict
and indifferent comparisons in both directions. Coordinatewise weak
improvement implies weak preference.
For every bet \(x\), there is a unique real \(C(x)\) such that
\[
                    x\sim C(x)\mathbf1,
       \qquad \min_i x_i\leq C(x)\leq\max_i x_i.
\]
This \emph{cash equivalent} satisfies
\[
\begin{aligned}
 C(t\mathbf1)&=t,& C(x+y)&=C(x)+C(y),\\
 C(\lambda x)&=\lambda C(x)&& (\lambda\in\mathbb R),\\
 |C(x)-C(y)|&\leq\|x-y\|_\infty.&
\end{aligned}
\]
Furthermore \(x\succeq y\) exactly when \(C(x)\geq C(y)\).
This assertion includes zero and negative scalars and the one-state case.
\end{hblemma}

\section{The probability representation}
\label{sec:p17-subjective-probability-representation}

\begin{hbdefinition}{def:p17-bet-probability-vector}{Probability vector and linear evaluation of a bet}
A \emph{probability vector} on \(S\) is a vector
\(\pi=(\pi_1,\ldots,\pi_K)\) with \(\pi_i\geq0\) and
\(\sum_i\pi_i=1\). Define
\[
                 \pi\cdot x\defeq\sum_{i=1}^K\pi_i x_i,
       \qquad \mathbb P_\pi(A)\defeq\sum_{i\in A}\pi_i
                                      \quad(A\subseteq S).
\]
The latter is the finite probability measure with those atomic masses.
The linear value \(\pi\cdot x\) is its expectation of the payoff
function \(i\mapsto x_i\), in the sense of
\cref{def:p15-finite-expectation}.
\end{hbdefinition}

\begin{hbtheorem}{thm:p17-finite-bet-representation}{proved=annex:p17.finite-bet-representation}{Finite additive betting representation}
A relation on \(\mathbb R^K\) is a complete, transitive preference
satisfying the three bet axioms if and only if there is a probability
vector \(\pi\) such that
\[
                   x\succeq y
                      \quad\Longleftrightarrow\quad
                   \pi\cdot x\geq\pi\cdot y.
\]
The representing probability vector is unique, and
\[
                        \pi_i=C(e_i),\qquad C(x)=\pi\cdot x.
\]
Zero coordinates of \(\pi\) are permitted. In the one-state case
the unique coordinate is one, and the order is the usual order of
the sure monetary payoff.
\end{hbtheorem}

The source proves its representation by separation of two convex sets
\cite[p.~103]{Rubinstein2015LectureNotes}.
The annex gives the complete cash-equivalent construction of the
separating probability vector and verifies the same two halfspaces.
This supplies existence without invoking an unstated separation theorem.
The result is the finite additive monetary-bet presentation; it does
not substitute for a general Savage representation with separately
identified consequence utility and subjective probability.

\begin{hbdefinition}{def:p17-bet-comparison-cones}{Acceptable and strictly negative comparison sets}
Put
\[
 \mathcal U\defeq\{x:x\succeq\mathbf0\},\qquad
 \mathcal D\defeq\{x:\mathbf0\succ x\}.
\]
A subset of \(\mathbb R^K\) is \emph{convex} if it contains
\(\alpha x+(1-\alpha)y\) for any two of its members and
\(0\leq\alpha\leq1\). A \emph{halfspace} with normal \(\pi\)
and threshold zero is specified by one linear inequality such as
\(\pi\cdot x\geq0\); its strict counterpart uses \(<0\).
\end{hbdefinition}

\begin{hbtheorem}{thm:p17-bet-halfspaces}{proved=annex:p17.bet-halfspaces}{The geometric form of the betting comparison}
Under the bet axioms, \(\mathcal U\) and \(\mathcal D\) are
nonempty disjoint convex sets with union \(\mathbb R^K\).
The first is closed and the second is open. For the unique probability
vector of the representation,
\[
          \mathcal U=\{x:\pi\cdot x\geq0\},\qquad
          \mathcal D=\{x:\pi\cdot x<0\}.
\]
The normal is nonzero, nonnegative coordinatewise and normalized to
sum to one. Both the boundary and its orientation are therefore fixed
by the actual comparisons of sure payoffs.
\end{hbtheorem}

\begin{hbtheorem}{thm:p17-bet-events-null-states}{proved=annex:p17.bet-events-null-states}{Event bets and null states}
For the indicator payoff \(\mathbf1_A=\sum_{i\in A}e_i\),
\[
                         C(\mathbf1_A)=\mathbb P_\pi(A).
\]
Changing the payoff only in state \(i\) leaves every bet indifferent
to its changed version exactly when \(\pi_i=0\).
The stronger condition that every coordinatewise improvement, strict
in at least one coordinate, is strictly preferred holds exactly when
\(\pi_i>0\) for every state.
\end{hbtheorem}

\begin{hbmethod}{def:p17-bet-elicitation}{Recover and use a finite betting model}
For the fixed state list and payoff units:
\begin{enumerate}
\item Specify the payoff of every contemplated act in every state, including
losses and the common settlement time. The states must be exhaustive and
mutually exclusive in this model.
\item State or assess the completeness, transitivity, continuity, monotonicity
and common-vector additivity assumptions. The theorem does not establish
that observed responses satisfy them.
\item Determine the sure-payment equivalents \(C(e_i)\) of the unit
state bets. Under the axioms these numbers are nonnegative and sum to one.
If observations do not have that property, the exact representation has
not been established by those observations.
\item Evaluate each candidate bet by \(\sum_iC(e_i)x_i\), and retain
every maximizer of the specified finite menu. Treat elicitation error and
uncertain state descriptions as model uncertainties, not algebraic identities.
\end{enumerate}
\end{hbmethod}

\begin{hbexample}{def:p17-bet-worked-example}{A stipulated three-state comparison}
Suppose \(C(e_1)=3/10\), \(C(e_2)=1/2\), and \(C(e_3)=1/5\).
They sum to one. The bet \(x=(10,0,-5)\) has
\[
                         C(x)=3-1=2,
\]
and is therefore preferred to a sure payment of one. The unit payoff
on \(A=\{1,3\}\) has cash equivalent \(3/10+1/5=1/2\).
These are hypothetical model inputs, not inferred frequencies.

Alternatively the vector \((1/2,1/2,0)\) satisfies all three axioms
through the representation, but changes in the third-state payoff have
no effect on its comparisons. The axioms permit a zero-probability
state; they do not assert that the state is logically impossible.
\end{hbexample}

\section{Finite fair prices and a sure-loss criterion}
\label{sec:p17-finite-price-coherence}

\begin{hbdefinition}{def:p17-finite-price-coherence}{Quoted event contracts and sure-loss coherence}
On the same finite nonempty state set, let \(q:2^S\to\mathbb R\)
be a quoted price for every unit event payoff. The net payoff after
paying that sure price is
\[
                       b_A\defeq\mathbf1_A-q(A)\mathbf1.
\]
A \emph{portfolio} is any finite real linear combination
\(x=\sum_{j=1}^m a_jb_{A_j}\), with \(m\geq0\).
Negative coefficients allow the reverse transaction at the same quote;
there are no transaction charges, quantity restrictions or differing
settlement dates in this mathematical model. The empty portfolio is zero.

Call the quote schedule \emph{coherent against a strict sure loss} if
no such portfolio has \(x_i<0\) in every state.
This is a finite contract criterion. It does not assert that an
individual actually accepts every transaction permitted by the model.
\index{Dutch book!finite event prices}\index{coherence!finite event prices}
\end{hbdefinition}

\begin{hbtheorem}{thm:p17-finite-price-coherence}{proved=annex:p17.finite-price-coherence}{Probability prices are exactly the coherent finite quotes}
The quote schedule is coherent against a strict sure loss if and only
if it is a probability measure on \(2^S\). Equivalently, there is
a unique probability vector \(\pi\) with
\[
                            q(A)=\sum_{i\in A}\pi_i.
\]
For probability quotes every portfolio has \(\pi\)-expected net
payoff zero. If all \(\pi_i>0\), there is also no portfolio which
is nonpositive in every state and negative in at least one. If a
coordinate \(\pi_i=0\), that latter stronger exclusion fails:
the portfolio \(-b_{\{i\}}=-e_i\) has exactly such a payoff.
\end{hbtheorem}

This theorem is an explicit finite-sum consequence proved here from
the stated contract model. It differs from the source's earlier
exchange-cycle discussion about the independence axiom
\cite[pp.~100--101]{Rubinstein2015LectureNotes}, whose steps include
conditional replacement and payments. A conclusion about that dynamic
exchange requires its own assumptions about consent, payments, timing
and the persistence of preferences. The finite portfolio theorem
does not supply those assumptions implicitly.

\section{The force of the additivity assumption}
\label{sec:p17-bet-additivity-boundary}

\begin{hbdefinition}{def:p17-nonlinear-bet-example}{A continuous nonlinear evaluation}
On \(\mathbb R^2\), define
\[
                         V(x)\defeq-e^{-x_1}-e^{-x_2},
\]
and compare bets by their \(V\)-values.
\end{hbdefinition}

\begin{hbtheorem}{thm:p17-nonlinear-bet-example}{proved=annex:p17.nonlinear-bet-example}{Beliefs alone do not imply additive monetary evaluation}
This preference is complete, transitive, continuous, and strictly
prefers an improvement in both states. It violates common-vector
additivity. Indeed, for \(x=(1,0)\), \(y=(0,1)\), and \(z=(0,1)\),
one has \(x\sim y\) but \(x+z\succ y+z\).
It therefore has no representation by a linear probability-weighted
average of monetary payoffs on the whole bet space.
\end{hbtheorem}

\begin{hbwarning}{def:p17-subjective-probability-boundary}{A representation is conditional on its decision model}
The counterexample can itself be evaluated as equal-weight expected
utility of the transformed payoffs \(-e^{-x_i}\), up to a positive
factor. Thus failure of additive monetary evaluation does not show
that an agent lacks probabilities or cannot use expected utility.
The finite betting theorem identifies precisely the stronger model
licensed by its own additivity condition.

Subjective probabilities obtained from a preference are model-relative
weights, not measurements of objective chance, guarantees of calibration,
or a proof of the correctness of the underlying state description.
The source explicitly distinguishes its small betting presentation
from Savage's larger axiomatization
\cite[p.~102]{Rubinstein2015LectureNotes}. General acts, nonlinear
consequence utility, ambiguity, updating and dynamic consistency
require their separate treatments.
\end{hbwarning}
